\section{Discussion}
\label{sec:discussion}

\subsection{Limitations}

\textbf{Single-node deployment.}
Minikube does not replicate the network topology of a production multi-node
cluster.
In particular, east-west traffic between nodes traverses an overlay network
that may differ from cloud provider CNI plugins (e.g., Cilium, Calico).

\textbf{Contamination assumption.}
The iForest contamination parameter is set globally at 5\%.
In practice, attack prevalence varies widely across time and environment.
An adaptive contamination estimator (e.g., via a sliding-window attack-rate
estimator fed from confirmed Falco alerts) would improve calibration.

\textbf{LSTM sequence length.}
We use a fixed window of $T = 10$ flows.
Slow exfiltration campaigns spanning hours would benefit from a longer
context, but this increases memory and inference latency.

\subsection{Future Work}

\begin{enumerate}
  \item \textbf{Online learning.} Update iForest and the AE continuously
        as new labelled alerts accumulate, reducing concept drift.
  \item \textbf{Graph-based features.} Encode pod-to-pod communication
        as a dynamic graph and add GNN-derived centrality features to the
        input vector.
  \item \textbf{Federated learning.} Train across multiple clusters without
        sharing raw flow data, preserving multi-tenancy isolation.
  \item \textbf{Explainability.} Integrate SHAP values to surface which
        features most contributed to each anomaly score, aiding analyst
        triage.
\end{enumerate}
